% !PRINTED 309
\weberpage{309}
%% [page furniture, not body text -- commented out 2026-08-10] Page furniture (running
%% head, folio, signature) is excluded from this edition by standing policy.
%% \noindent \S.\,80. \hfill Gleichungen 8ten Grades.

%<<CONT
so dass $R$ die gesammte ternäre lineare Congruenzgruppe für den Modul 2 ist.

Wir betrachten ein System von acht unabhängigen Veränderlichen $X_{z_1,z_2,z_3}$, worin die Indices $z_1, z_2, z_3$ nach dem Modul 2 zu nehmen sind.

Hieraus bilden wir die sieben Functionen (\gesperrt{Lagrange'sche} Resolventen, Bd.\,I, \S.\,164).
\weq{1}{\overset{z}{\Sigma} (-1)^{\Sigma z\xi} X_{z_1,z_2,z_3} = \Psi_{\xi_1,\xi_2,\xi_3},}
worin zur Abkürzung
\[
\Sigma z\xi = z_1\xi_1 + z_2\xi_2 + z_3\xi_3
\]
gesetzt ist, und $\xi_1, \xi_2, \xi_3$ je ein volles Restsystem nach dem Modul 2, jedoch mit Ausschluss der Combination $0, 0, 0$ durchlaufen.

Wenden wir auf die Indices $z_i$ der Grössen $X$ die Substitutionen der Gruppe $Q$ an, so erleiden die $X$ die Permutationen einer Gruppe, die wir gleichfalls mit $Q$ bezeichnen. Die Aenderungen der $\Psi$, die diesen Permutationen entsprechen, ergeben sich aus (1), wenn wir auf die Indices der $X$ die Substitution $(z_i, z_i+h)$ anwenden:
\weq{2}{\Sigma (-1)^{\Sigma z\xi} X_{z_1+h_1, z_2+h_2, z_3+h_3} = \overset{z}{\Sigma} (-1)^{\Sigma(z+h)\xi} X_{z_1,z_2,z_3} = (-1)^{\Sigma h\xi} \Psi_{\xi_1,\xi_2,\xi_3}.}

\gesperrt{1. Die Functionen $\Psi$ ändern also durch die Permutationen von $Q$ nur ihr Zeichen, und die Quadrate der $\Psi$ bleiben durch $Q$ ungeändert.}

Es ist ferner der Einfluss einer Substitution $\sigma$ aus $S$ auf die Functionen $\Psi$ festzustellen.

Bezeichnen wir zu diesem Zwecke mit $\varrho$ die \gesperrt{transponirte Substitution} zu $\sigma$ und setzen
\weq{3}{z' = \sigma(z), \quad \xi = \varrho(\xi'), \quad \xi' = \varrho^{-1}(\xi),}
so ist (\S.\,37, 9.):
\weq{4}{\Sigma \xi z = \Sigma \xi' z'.}

Machen wir die Substitutionen $\sigma$ der Gruppe $S$ in den Indices von $X$, so erhalten wir eine mit $S$ zu bezeichnende Permutationsgruppe der $X$, und aus (1) ergiebt sich
\[
\overset{z}{\Sigma} (-1)^{\Sigma z\xi} X_{z_1', z_2', z_3'} = \overset{z}{\Sigma} (-1)^{\Sigma \xi' z'} X_{z_1', z_2', z_3'},
\]
%CONT>>
